\section{A Ética da Simulação Algorítmica}
\label{sec:etica-simulacao}

\subsection{O Imperativo da Transparência Algorítmica (XAI)}

A predição simulada traz consigo a opacidade matemática. Se um algoritmo de IA determina que a extração de pré-molares é desnecessária em um caso limítrofe de ortodontia, o clínico deve ter a capacidade técnica de auditar como essa decisão foi fundamentada. O Regulamento Europeu de Inteligência Artificial (\textit{EU AI Act}) classifica os softwares de IA médica como sistemas de alto risco, exigindo que sejam providos de \textit{Explainable AI} (IA Explicável). Na odontologia, isso implica em fornecer ao dentista não apenas o plano A, mas mapas de calor (\textit{saliency maps}) ou escores de incerteza que permitam avaliar o peso que o algoritmo atribuiu ao biotipo gengival ou à densidade óssea.

\subsection{Viés Algorítmico e \textit{Data Bias} Estrutural}

Modelos preditivos treinados exclusivamente em bancos de dados de pacientes europeus caucasianos tenderão a errar substancialmente quando aplicados à geometria craniofacial da população miscigenada brasileira. Um algoritmo que não foi exposto durante a sua fase de treinamento (Deep Learning) a pacientes com histórico de doença periodontal severa (comum na população brasileira de baixa renda) pode subestimar o risco de falha em um planejamento de implantes. O viés estrutural pode, assim, institucionalizar o erro médico por razões demográficas. A correção ética exige a criação de \textit{datasets} representativos da diversidade nacional e a auditoria contínua de equidade do algoritmo \cite{roy2025editorial}.
